\section*{Chapter \ref{ch:b2:alpine-orogen} --- Anatomy of a Mountain Belt: The Alps}

\begin{solution}{exo:b2:alpine-orogen:1}
(a) $e = 1 - 27.0/36.0 = 0.25$, that is \qty{25}{\percent}; $S = 36.0 - 27.0 =
\qty{9.0}{km}$. (b) The bed must have kept its length (no stretching or thinning along it),
the section must be parallel to the transport direction, and the pin lines must mark places
where the layers did not slip past each other.
\end{solution}

\begin{solution}{exo:b2:alpine-orogen:2}
(a) A \omterm{def:b2:alpine-orogen:klippe}{klippe}: the gneiss is an isolated remnant of a \omterm{def:b2:alpine-orogen:nappes}{nappe} resting on its autochthon;
erosion has removed the rest of the sheet all round it. (b) A \omterm{def:b2:alpine-orogen:klippe}{tectonic window}: erosion has
cut through the gneiss \omterm{def:b2:alpine-orogen:nappes}{nappe} and exposed the younger rocks beneath its basal thrust, which
therefore dips away from the window on every side.
\end{solution}

\begin{solution}{exo:b2:alpine-orogen:3}
(a) Ultramafic complex (the upper mantle); gabbroic complex (the lower oceanic crust);
sheeted \omterm{def:b2:magma-emplacement:bodies}{dyke} complex (the middle crust); volcanic complex, pillowed (the upper crust).
(b) The sheeted \omterm{def:b2:magma-emplacement:bodies}{dyke} complex. (c) Most crust of the wide oceans is consumed by subduction;
the \omterm{def:b2:alpine-orogen:ophiolite}{ophiolites} that survive on land formed mostly in small ocean basins or between island
arcs, the Tethyan ones where a continent thinned until a small ocean opened.
\end{solution}

\begin{solution}{exo:b2:alpine-orogen:4}
(a) Monte Rosa and Grand St Bernard \omterm{def:b2:alpine-orogen:nappes}{nappes}: the European side (the Briançonnais and the
European passive margin); Zermatt--Saas and Combin zone: the Piemont--Liguria ocean; Dent
Blanche system: the margin of Apulia. (b) Apulia, which ends on top. (c) The upper plate is
not subducted, and the slices scraped off the lower plate stack beneath it in the order the
lower plate arrived: far ocean first, near margin last and lowest. Later folding of the pile,
or out-of-sequence thrusts, can upset the order.
\end{solution}

\begin{solution}{exo:b2:alpine-orogen:5}
(a) $2\pi G\,\Delta\rho = 2\pi\times6.674\times10^{-11}\times500 =
\qty{2.097e-7}{s^{-2}}$; $h = 1.70\times10^{-3}/2.097\times10^{-7} = \qty{8.1}{km}$.
(b) $2.097\times10^{-7}\times25\,000 = \qty{5.24e-3}{m.s^{-2}}$, about
\qty{-524}{mGal}. (c) The root is not an infinite slab: it is a body some tens of kilometres
deep and only about as wide as the belt, so it pulls less on a point at the surface than a
slab of the same thickness extending to infinity would; dense bodies such as the Ivrea zone
add positive anomalies, and the density contrast at depth may be smaller than assumed.
\end{solution}

\begin{solution}{exo:b2:alpine-orogen:6}
(a) $164 - 44 = \qty{120}{Myr}$. (b) $z = P/(\bar\rho g)$:
$2.8\times10^9/(3000\times9.8) = \qty{95.2}{km}$; $2.8\times10^9/(3300\times9.8) =
\qty{86.6}{km}$. (c) ``More than \qty{90}{km}'' lies inside this range: the depth depends on
the mean density chosen for the overburden as much as on the pressure.
\end{solution}

\begin{solution}{exo:b2:alpine-orogen:7}
(a) A limb of length $L$ dipping at $\theta$ spans $L\cos\theta$ horizontally; all limbs
have the same dip, so $l = l_0\cos\theta$ and $e = 1 - l/l_0 = 1 - \cos\theta$.
(b) $\cos\theta = 0.70$: $\theta = \qty{45.6}{\degree}$. (c) $1 - \cos\qty{40}{\degree} =
\qty{23.4}{\percent}$, as in the chevron figure; $1 - \cos\qty{60}{\degree}
= \qty{50}{\percent}$.
\end{solution}

\begin{solution}{exo:b2:alpine-orogen:8}
(a) $l_0 = 118/2.0 = \qty{59.0}{km}$; $S = 59.0 - 42 = \qty{17.0}{km}$; $e = 17.0/59.0 =
\qty{28.8}{\percent}$. (b) $A_e = 118 - 42\times2.0 = \qty{34.0}{km^2}$; $d = 34.0/17.0 =
\qty{2.0}{km}$. (c) Yes: the \omterm{def:b2:alpine-orogen:decollement}{décollement} lies at the base of the \qty{2.0}{km} cover, below
its top; the cover is detached on its base.
\end{solution}

\begin{solution}{exo:b2:alpine-orogen:9}
(a) The Airy root is $4.5h$: at \qty{46.5}{\degree} N, $4.5\times2.74 = \qty{12.3}{km}$, a
Moho at \qty{42.3}{km}; at \qty{45.55}{\degree} N, $4.5\times0.65 = \qty{2.9}{km}$, a Moho
at \qty{32.9}{km}. (b) Under the high ground the Moho is \qty{6.0}{km} shallower than Airy
predicts; under the edge of the Po plain it is \qty{28.3}{km} deeper. (c) The root is not
compensated locally: it is displaced towards the south and larger than the relief requires.
The plate bends as a whole under the loads on it (the regional support of a Year 3
chapter), and one of those loads, the Ivrea body, is not topography at all.
\end{solution}

\begin{solution}{exo:b2:alpine-orogen:10}
(a) $5.1/11.3 = 0.45$: less than half. (b) $5.1 + 6.2 = \qty{11.3}{km}$. The top of the
Jurassic lies above the upper detachment; the \qty{6.2}{km} carried along that detachment
passed beneath the folds of the section and was spent in structures further north, so the
upper bed never shortened by that amount within the section. (c) A single marker bed records
only the shortening of the layers it belongs to; every detachment above or below it can
carry displacement in or out of the section. Balance several beds, and check with areas.
\end{solution}

\begin{solution}{exo:b2:alpine-orogen:11}
(a) For an offset of \qty{100}{km}, $\sqrt{200^2 + 100^2} = \qty{223.6}{km}$; for
\qty{50}{km}, $\sqrt{200^2 + 50^2} = \qty{206.2}{km}$; for \qty{150}{km},
$\sqrt{200^2 + 150^2} = \qty{250.0}{km}$. (b) $\arctan(100/200) = \qty{26.6}{\degree}$ from
north--south, a convergence line close to the N\qty{150}{\degree}E quoted by Rosenberg and
colleagues. (c) The offset is uncertain by $\pm\qty{50}{km}$, which moves the bulk
convergence by about \qty{20}{km} either way; the third figure of the result is far below
that uncertainty.
\end{solution}

\begin{solution}{exo:b2:alpine-orogen:12}
(a) $210/12.9 = \qty{16.3}{Myr}$ to $210/8.5 = \qty{24.7}{Myr}$ (\unit{km} divided by
\unit{mm.yr^{-1}} gives millions of years). (b) More than $4\times60 = \qty{240}{km}$; at
least $240 - 60 = \qty{180}{km}$ of the basin is now under the Alpine thrusts or has been
eroded. (c) The basin forms in front of the thrust wedge and migrates with it: its inner
part is always being overridden by the thrusts that load it, and its oldest fill is caught
in the thrusts, as the Subalpine \omterm{def:b2:alpine-orogen:foreland}{Molasse} is.
\end{solution}

\begin{solution}{pb:b2:alpine-orogen:1}
\noindent\textbf{Part I.}

\textbf{1.} $3.0/\cos\qty{35}{\degree} = \qty{3.66}{km}$; $2.5/\cos\qty{50}{\degree} =
\qty{3.89}{km}$; $4.0/\cos\qty{40}{\degree} = \qty{5.22}{km}$.

\textbf{2.} $8.00 + 2\times3.66 + 5.00 + 2\times3.89 + 6.0/\cos\qty{10}{\degree} +
2\times5.22 + 10.0/\cos\qty{5}{\degree} = 8.00 + 7.32 + 5.00 + 7.78 + 6.09 + 10.44 + 10.04 =
\qty{54.67}{km}$; with the repeated stretch, $l_0 = \qty{60.7}{km}$.

\textbf{3.} The bed beneath the thrust sheet is the same bed, carried under the sheet by the
thrust: before thrusting it lay in line with the rest. Leaving it out would restore the
bed too short.

\textbf{4.} $e = 1 - 48.0/60.7 = 0.209$, that is \qty{20.9}{\percent}; $S = 60.7 - 48.0 =
\qty{12.7}{km}$.

\textbf{5.} $60.7 - 60.2 = \qty{0.5}{km}$, $0.5/60.7 = \qty{0.8}{\percent}$: within the
precision of the measurement; the section is acceptable.

\textbf{6.} $60.7 - 54.0 = \qty{6.7}{km}$, \qty{11}{\percent}: the two beds of one sheet
cannot restore to such different lengths. The section would not be admissible as drawn: a
structure missing or misdrawn, a pin line in the wrong place, movement out of the plane of
the section, or a detachment between the two beds.

\noindent\textbf{Part II.}

\textbf{7.} $146/2.4 = \qty{60.8}{km}$.

\textbf{8.} It differs from \qty{60.7}{km} by \qty{0.1}{km}: line length and area agree.

\textbf{9.} $A_e = 146 - 48.0\times2.4 = 146 - 115.2 = \qty{30.8}{km^2}$.

\textbf{10.} $d = 30.8/12.7 = \qty{2.43}{km}$, equal to the \qty{2.4}{km} of the cover to
the precision of the data: the \omterm{def:b2:alpine-orogen:decollement}{décollement} is the base of the cover.

\textbf{11.} $160/2.4 = \qty{66.7}{km}$, \qty{6.0}{km} more than the line length. Thickening
of the evaporites at the base of the cover, flowing into the cores of anticlines, adds area
without adding length; area balance with a constant thickness then fails.

\textbf{12.} The cover moved on a \omterm{def:b2:alpine-orogen:decollement}{décollement} at its base and the basement was not
involved. Beneath the belt the top of the basement should run flat or dip gently and
evenly, unbroken by the thrusts above it.

\noindent\textbf{Part III.}

\textbf{13.} $22 + 18 + 25 = \qty{65}{km}$; $S = 65 - 15 = \qty{50}{km}$.

\textbf{14.} $e = 50/65 = \qty{76.9}{\percent}$, against \qty{20.9}{\percent} in the
foreland belt: the \omterm{def:b2:alpine-orogen:nappes}{nappes} are far more shortened.

\textbf{15.} $30 - 22 = \qty{8}{km}$ ($e = \qty{26.7}{\percent}$). The basement itself is
shortened: \omterm{def:b2:alpine-orogen:skin}{thick-skinned tectonics}.

\textbf{16.} Foreland belt $12.7/(11 - 4) = \qty{1.81}{km/Myr}$, that is
\qty{1.81}{mm} a year; \omterm{def:b2:alpine-orogen:nappes}{nappes} $50/(30 - 20) = \qty{5.0}{mm}$ a year.

\textbf{17.} The \omterm{def:b2:alpine-orogen:nappes}{nappes} (30 to \qty{20}{Ma}) deformed before the foreland belt (11 to
\qty{4}{Ma}): deformation migrated outwards, towards the foreland. This is the order of the
Alps: the Glarus thrust in the Oligocene and Miocene, the external massifs in the Miocene,
the Jura in the late Miocene and Pliocene.

\textbf{18.} $12.7 + 50 + 8 = \qty{70.7}{km}$.

\noindent\textbf{Part IV.}

\textbf{19.} $\qty{0.5}{cm.yr^{-1}}\times\qty{40}{Myr} = 0.005\times4\times10^7 =
\qty{2.0e5}{m} = \qty{200}{km}$.

\textbf{20.} $200 + 18 = \qty{218}{km}$.

\textbf{21.} Only what is preserved or inferred is counted: parts of the ocean may have
been subducted before \qty{85}{Ma} or faster than the mean rate, \omterm{def:b2:alpine-orogen:ophiolite}{ophiolite} slices may have
been eroded or carried out of the section, and convergence oblique to the section brings
in ocean floor the section does not see.

\textbf{22.} Internal zone $60 - 20 = \qty{40}{km}$ ($e = \qty{66.7}{\percent}$); retro-wedge
$70 - 48 = \qty{22}{km}$ ($e = \qty{31.4}{\percent}$).

\textbf{23.} $12.7 + 50 + 8 + 40 + 22 = \qty{132.7}{km}$; within the range of the estimates
for the Central Alps, about 125 to \qty{160}{km} of collisional shortening (before adding
strike slip).

\textbf{24.} Restored width $153 + 132.7 = \qty{285.7}{km}$; restored crust
$285.7\times30 = \qty{8571}{km^2}$; present crust $153\times45 = \qty{6885}{km^2}$. Missing:
$8571 - 6885 = \qty{1686}{km^2}$, \qty{19.7}{\percent} of the restored crust; it has been
eroded (and shed into the \omterm{def:b2:alpine-orogen:foreland}{foreland basins}) or subducted with the lower crust into the
mantle. Without those losses the crust would be $8571/153 = \qty{56}{km}$ thick on average.
The total collisional shortening across the section is
\textbf{\qty[text-series-to-math=true]{133}{km}}.
\end{solution}
